\section{Sensitivity and query complexity}

For total Boolean functions on $n$ bits, write $s$ for maximum sensitivity,
$bs$ for maximum block sensitivity, $C$ for certificate complexity, and $D$
for deterministic decision-tree depth. A subscript $b$ restricts a maximum
to inputs with output $b$; it is zero if there are no such inputs.

\begin{result}
\subsection{Polynomial relations between query measures}
For every total Boolean function $f$,
\[
C(f)\le s(f)\,bs(f)\le bs(f)^2,\qquad
D(f)\le C(f)\,bs(f)\le bs(f)^3.
\]
The formal development also gives $C(f)\le2s(f)^5$ and $D(f)\le4s(f)^9$,
using the existing sensitivity theorem and degree bounds.
These are upper bounds for exact deterministic queries, with no promise
on the input.\citep{Nisan1991,BuhrmanDeWolf2002,Huang2019,NisanSzegedy1994}
\decl{Complexity/Computability/QueryComplexity/Relations.lean}{200}{Cslib.QueryComplexity.certificateComplexity_le_sensitivity_mul_blockSensitivity}
\decl{Complexity/Computability/QueryComplexity/Relations.lean}{206}{Cslib.QueryComplexity.certificateComplexity_le_blockSensitivity_sq}
\decl{Complexity/Computability/QueryComplexity/Relations.lean}{366}{Cslib.DecisionTree.complexity_le_certificateComplexity_mul_blockSensitivity}
\decl{Complexity/Computability/QueryComplexity/Relations.lean}{382}{Cslib.DecisionTree.complexity_le_blockSensitivity_pow_three}
\decl{Complexity/Computability/QueryComplexity/SensitivityBounds.lean}{50}{Cslib.QueryComplexity.certificateComplexity_le_two_mul_sensitivity_pow_five}
\decl{Complexity/Computability/QueryComplexity/SensitivityBounds.lean}{64}{Cslib.DecisionTree.complexity_le_four_mul_sensitivity_pow_nine}
\end{result}

\begin{result}
\subsection{Monotone measures and symmetric evasiveness}
For monotone $f$, $C(f)=bs(f)=s(f)$ and $D(f)\le s_0(f)s_1(f)\le s(f)^2$.
Separately, every nonconstant symmetric Boolean function on $n$ bits is
evasive: $D(f)=n$. Symmetry means dependence only on Hamming weight.
Thus, for $\mathrm{Th}^n_j$, which accepts iff at least $j$ bits are true,
with $1\le j\le n$,
\[
D(\mathrm{Th}^n_j)=n,\qquad
C(\mathrm{Th}^n_j)=bs(\mathrm{Th}^n_j)=s(\mathrm{Th}^n_j)
 =\max(j,n+1-j).
\]
The equality of the three local-change and certificate measures requires
monotonicity; it is not asserted for arbitrary functions.
\citep{Nisan1991,BuhrmanDeWolf2002}
\decl{Complexity/Computability/QueryComplexity/Monotone.lean}{128}{Cslib.QueryComplexity.certificateComplexity_eq_sensitivity_of_monotone}
\decl{Complexity/Computability/QueryComplexity/Monotone.lean}{136}{Cslib.QueryComplexity.blockSensitivity_eq_sensitivity_of_monotone}
\decl{Complexity/Computability/QueryComplexity/Monotone.lean}{150}{Cslib.DecisionTree.complexity_le_sensitivityFor_mul_sensitivityFor_of_monotone}
\decl{Complexity/Computability/DecisionTree/Symmetric.lean}{160}{Cslib.DecisionTree.complexity_eq_card_of_isSymmetric}
\decl{Complexity/Computability/QueryComplexity/Threshold.lean}{140}{Cslib.BitString.sensitivity_threshold}
\end{result}

\begin{result}
\subsection{Approximate degree of disjunction}
Let $\widetilde{\deg}_{1/3}(f)$ be the least degree of a real polynomial
approximating the $0/1$ value of $f$ to error at most $1/3$ at every Boolean input.
For disjunction on $n$ bits,
\[
n\le6\,\widetilde{\deg}_{1/3}(\mathrm{OR}_n)^2,\qquad
\widetilde{\deg}_{1/3}(\mathrm{OR}_n)\le\lfloor\sqrt{2n}\rfloor+1.
\]
These explicit finite bounds establish the square-root order of growth.
The upper bound uses a Chebyshev polynomial; no optimal leading constant is claimed.
\citep{NisanSzegedy1994,BuhrmanDeWolf2002}
\decl{Complexity/Computability/QueryComplexity/ApproximateDegree/Or.lean}{85}{Cslib.QueryComplexity.card_le_six_mul_approxDegree_threshold_one_right_sq}
\decl{Complexity/Computability/QueryComplexity/ApproximateDegree/Or.lean}{168}{Cslib.QueryComplexity.approxDegree_threshold_one_right_le_sqrt}
\end{result}

\begin{result}
\subsection{Superquadratic separation of sensitivity and block sensitivity}
There is a fixed real exponent $\alpha>2$ and a family of nonconstant Boolean
functions $F_m$ on positive, varying arities such that $F_m(0)=0$ and
\[
s(F_m)^\alpha\le bs(F_m,0),\qquad bs(F_m,0)\longrightarrow\infty.
\]
Here $0$ is the all-false input, and $bs(f,0)$ is block sensitivity at that
specific input. In particular, for every real $c$ there exists a nonconstant
$f$ with $c\,s(f)^2<bs(f)$, so no universal quadratic upper bound holds.
The exponent is existential, not an assertion for every $\alpha>2$.
The linked construction ports OpenAI's formalization through the earlier
complexitylib development and connects it to the library's public measures.
\citep{OpenAI2026Sensitivity,Schlesinger2026Sensitivity}
\decl{Complexity/Computability/QueryComplexity/SensitivitySeparation.lean}{108}{Cslib.QueryComplexity.exists_sensitivity_rpow_le_blockSensitivityAt}
\decl{Complexity/Computability/QueryComplexity/SensitivitySeparation.lean}{99}{Cslib.QueryComplexity.exists_mul_sensitivity_sq_lt_blockSensitivity}
\end{result}

\begin{result}
\subsection{Element distinctness for entries below \texorpdfstring{$n$}{n} on the word RAM}
On the bounded-word RAM, whose registers and memory cells hold $w$-bit words, whose every
instruction costs one unit, and on which writing a constant or an arithmetic result of more than
$w$ bits faults, a single fixed program decides whether a list of $n$ naturals, each smaller than
$n$, has no duplicates: at every word size
$w\ge 2(\lfloor\log_2\max(1,n)\rfloor+1)$ it halts within
$16(n+1)$ instructions, the final halt included, with output $1$ exactly when the entries are
distinct. Every promised input fits at this threshold, including the empty list at $w=2$.
Hence the problem is in linear time on the word RAM. The program allocates a zeroed
table of $n$ flags and uses each entry as an address into it (direct addressing), so it is not a
comparison algorithm. The bound on the entries is part of the instance type; unrestricted element
distinctness is not treated. Uniform decision requires one program for all word sizes above the
threshold, as in the transdichotomous model. The machine, its time classes and its structured
programs, which compile with identical instruction counts, are adapted from the Lean
formalization accompanying the tromino undecidability paper.
\citep[Sections 2.2 and 11.1]{CLRS2009}\citep{FredmanWillard1993,Hagerup1998}\citep{Tromino2025}
\declat{4d59a90c037f3ae099131540a2da2d981e97a003}{Complexity/Algorithms/WordRAM/ElementDistinctness.lean}{287}{Cslib.Algorithms.WordRAM.ElementDistinctness.decidesInTime_nodup}
\declat{4d59a90c037f3ae099131540a2da2d981e97a003}{Complexity/Algorithms/WordRAM/ElementDistinctness.lean}{296}{Cslib.Algorithms.WordRAM.ElementDistinctness.linearTime_nodup}
\declat{4d59a90c037f3ae099131540a2da2d981e97a003}{Complexity/Algorithms/WordRAM/ElementDistinctness.lean}{148}{Cslib.Algorithms.WordRAM.ElementDistinctness.inputFits}
\declat{4d59a90c037f3ae099131540a2da2d981e97a003}{Complexity/Computability/WordRAM/Time.lean}{132}{Cslib.WordRAM.Program.DecidesInTime}
\declat{4d59a90c037f3ae099131540a2da2d981e97a003}{Complexity/Computability/WordRAM/Compile.lean}{135}{Cslib.WordRAM.Exec.runs}
\end{result}

\begin{result}
\subsection{Additive word-RAM simulation with charged preparation}
Let $p$ be a fixed word-RAM program with $r$ named registers, using only constants, copies,
addition, truncated subtraction, indirect loads and stores, equality and unsigned comparison,
jumps and halting. Set
\[
R=r+5,\qquad \ell(m)=\lfloor\log_2\max(1,m)\rfloor+1,\qquad
A_r=8(\ell(R)+1),\quad B_r=(24+R)(\ell(R+5)+2).
\]
There is an executable compiler, independent of input and word size, into the library's
natural-number RAM with logarithmic instruction costs. If the input list $u$ has length and
entries below $2^w$, and $p$ halts on $u$ with output $v$ after $t$ source instructions,
including its final halt, the prepared target program halts with the same output within cost
\[
(A_r+B_r)(t+|u|+1)(w+1).
\]
Its external input is five zeros followed by $|u|$ and the entries of $u$. Installing this
fixed length-bearing convention is uncharged. The target then copies the header and words
backwards to address $R$, clears the prefix, initializes the compiler registers and executes
the translated program; every instruction, including the final output copy and jump, is charged.
The separate padded-input theorem starts with $R$ supplied zeros followed by $|u|$ and $u$,
and bounds the compiler execution by $A_r(t+1)(w+1)$.
Neither statement certifies an encoding reduction from an arbitrary bare list:
bare RAM inputs $[]$ and $[0]$ initialize identical states.

This is a forward simulation of successful runs, not reflection of source overflow or
falloff faults. Word-width access, multiplication, division, bitwise operations and shifts are
outside the supported fragment. The source model adapts the lean-trominoes formalization;
the target is the library's documented natural-number adaptation of Cook--Reckhow RAM.
No Turing-machine simulation or complexity-class equivalence is asserted.
\citep{CookReckhow1973,Tromino2025}
\declat{80a79abcb959601b9a3d570d9c6d3ad347e9f0fe}{Complexity/Computability/WordRAM/RAM.lean}{61}{Cslib.WordRAM.Instr.IsAdditive}
\declat{aec70b38fb3158a09850dc3640fba31e8e312f1e}{Complexity/Computability/WordRAM/RAMInput.lean}{55}{Cslib.WordRAM.ToRAM.preparedProgram}
\declat{aec70b38fb3158a09850dc3640fba31e8e312f1e}{Complexity/Computability/WordRAM/RAMInput.lean}{93}{Cslib.WordRAM.ToRAM.haltsWithResultInTime_prepared}
\declat{aec70b38fb3158a09850dc3640fba31e8e312f1e}{Complexity/Computability/WordRAM/RAMInput.lean}{61}{Cslib.WordRAM.ToRAM.haltsWithResultInTime_prepared_of_bitLength_le}
\declat{80a79abcb959601b9a3d570d9c6d3ad347e9f0fe}{Complexity/Computability/WordRAM/RAM.lean}{569}{Cslib.WordRAM.ToRAM.haltsWithResultInTime_padded}
\end{result}

\begin{result}
\subsection{Bounded-universe element distinctness with logarithmic RAM cost}
A single fixed program on the library's natural-number RAM decides whether a list of $n$
naturals, each below $n$, has no duplicates. It starts with five supplied zero scratch words,
a length header and the input words, and halts with output $1$ for a distinct list and $0$
otherwise within logarithmic instruction cost
\[
2735(n+1)\bigl(\lfloor\log_2(n+1)\rfloor+1\bigr).
\]
The program combines input preparation with the compiled twelve-register direct-addressing
word-RAM algorithm above. Every relocation, clearing, initialization, simulation and output instruction
is charged: preparation contributes at most $287(n+1)(\lfloor\log_2(n+1)\rfloor+1)$ and
compiled execution at most $2448(n+1)(\lfloor\log_2(n+1)\rfloor+1)$. The latter bound also
holds when the seventeen-zero padded workspace is supplied directly. The constants are not
optimized, and the statements include the empty input.

The external five-zero prefix, length header and input words are preloaded without charge.
The entry bound is a promise; unrestricted element distinctness and a reduction from arbitrary
encoded inputs are not established. This is a logarithmic-cost RAM result, with no claimed
Turing-machine time bound. Mathematical and model sources are direct addressing, Cook--Reckhow
RAM, and the lean-trominoes word-RAM formalization.
\citep[Section 11.1]{CLRS2009}\citep{CookReckhow1973,Tromino2025}
\declat{aec70b38fb3158a09850dc3640fba31e8e312f1e}{Complexity/Algorithms/WordRAM/ElementDistinctnessRAM.lean}{67}{Cslib.Algorithms.WordRAM.ElementDistinctness.ramProgram}
\declat{aec70b38fb3158a09850dc3640fba31e8e312f1e}{Complexity/Algorithms/WordRAM/ElementDistinctnessRAM.lean}{111}{Cslib.Algorithms.WordRAM.ElementDistinctness.haltsWithResultInTime_nodup}
\declat{aec70b38fb3158a09850dc3640fba31e8e312f1e}{Complexity/Algorithms/WordRAM/ElementDistinctnessRAM.lean}{99}{Cslib.Algorithms.WordRAM.ElementDistinctness.haltsWithResultInTime_nodup_padded}
\end{result}
